\section{Rational curves, cotangent positivity, and covers}
\label{sec:classification}

We first combine smooth canonical nonvanishing with classical results on
rational curves and cotangent tensors. We then use smooth abundance to
remove the abundance premise from the quasi-projective-cover results of
Claudon, H\"oring and Koll\'ar. Throughout this section,
\(\kappa(X)=\kappa(X,K_X)\). We retain the convention
\(\kappa(\mathrm{pt})=0\), and regard a point as rationally connected
but not uniruled.

\begin{corollary}[Uniruledness and Mumford's criterion]
\label{cor:uniruled-mumford}
Let \(X\) be a smooth connected projective complex variety.
\begin{enumerate}[label=\textup{(\roman*)}]
\item \(X\) is uniruled if and only if \(\kappa(X)=-\infty\).
\item \(X\) is rationally connected if and only if
\[
 H^0\bigl(X,(\Omega_X^1)^{\otimes m}\bigr)=0
 \quad\text{for every integer }m>0.
\]
\end{enumerate}
\end{corollary}
\begin{proof}
For positive-dimensional \(X\), the theorem of
Boucksom--Demailly--P{\u a}un--Peternell
\cite[Corollary~0.3 and Theorem~2.6]{BDPP2013} says that \(X\) is
non-uniruled exactly when \(K_X\) is pseudo-effective.
Corollary~\ref{cor:smooth-good-model} then gives \(\kappa(X)\geq0\):
the semiample canonical divisor on the good model has a nonzero section
in a positive multiple, and the pluricanonical section spaces agree.
Conversely, a nonzero pluricanonical section makes \(K_X\) pseudo-effective. This proves
\textup{(i)}.

For \textup{(ii)}, rational connectedness forces all positive covariant
tensor differentials to vanish, as recalled by Lazi{\'c} and Peternell.
For the converse, their Proposition~2.4 reduces Mumford's criterion in
dimension \(n\) to tensor nonvanishing for smooth projective varieties
of positive dimension at most \(n\) whose canonical class is
pseudo-effective \cite[Proposition~2.4]{LazicPeternell2017}.
If \(Y\) is such a variety of dimension \(d\), the preceding good-model
argument supplies a nonzero section of \(mK_Y\) for some integer \(m>0\).
Antisymmetrization embeds
\(\Omega_Y^d\) into \((\Omega_Y^1)^{\otimes d}\) over \(\C\);
tensoring this injection \(m\) times gives the required nonzero tensor
differential. This is the observation in
\cite[Remark~2.5]{LazicPeternell2017}, so their reduction applies in
every dimension. Both assertions for a point follow from our conventions.
\end{proof}

For a smooth connected projective complex variety \(X\), let \(R(X)\)
be a smooth projective model of the base of its maximal rationally
connected (MRC) fibration. This base is determined up to birational
equivalence and is non-uniruled unless it is a point. We use
\(\Omega_X^0=\OO_X\).

\begin{corollary}[Pluriforms and the rational quotient]
\label{cor:pluriforms}
\begin{enumerate}[label=\textup{(\roman*)}]
\item For every smooth connected projective complex variety \(X\),
\[
 \dim R(X)=
 \max_{\substack{p\in\Z_{\geq0},\,k\in\Z_{>0}\\
 H^0(X,\operatorname{Sym}^k(\Omega_X^p))\ne0}} p.
\]
\item A compact connected K\"ahler manifold \(X\) is rationally
connected if and only if
\[
 H^0\bigl(X,\operatorname{Sym}^k(\Omega_X^p)\bigr)=0
 \quad\text{for every pair of integers }p,k>0.
\]
\end{enumerate}
\end{corollary}
\begin{proof}
Write \(r=\dim R(X)\) in \textup{(i)}. Resolve the MRC map to a morphism
\(f:X'\to R(X)\), with \(X'\) smooth and projective. For a smooth
general fiber \(F\), the cotangent sequence is
\[
 0\longrightarrow\OO_F^{\oplus r}
 \longrightarrow\Omega_{X'}^1|_F
 \longrightarrow\Omega_F^1\longrightarrow0.
\]
If \(p>r\), every graded piece of the induced exterior-power filtration
on \(\Omega_{X'}^p|_F\) contains a positive exterior power of
\(\Omega_F^1\): at most \(r\) factors can come from
\(\OO_F^{\oplus r}\). Consequently, for \(k>0\), the graded pieces
of the induced filtration on
\(\operatorname{Sym}^k(\Omega_{X'}^p)|_F\) are direct summands of
sums of positive tensor powers of \(\Omega_F^1\). Here we use
antisymmetrization and symmetrization over \(\C\). Rational
connectedness of \(F\) makes all their sections vanish, by
Corollary~\ref{cor:uniruled-mumford}. A global section therefore
vanishes on every sufficiently general fiber, hence on \(X'\).
Birational invariance of global covariant tensor differentials on smooth
projective varieties gives the same vanishing on \(X\). This is the
upper-bound argument of \cite[Section~3]{BrunebarbeCampana2015}.

If \(r>0\), the smooth non-uniruled base has \(\kappa(R(X))\geq0\) by
Corollary~\ref{cor:uniruled-mumford}. A nonzero section of \(kK_{R(X)}\)
for some \(k>0\) pulls back to a nonzero section of
\(\operatorname{Sym}^k(\Omega_X^r)\). The pullback is taken on \(X'\)
and descends to \(X\) by the same birational invariance. If \(r=0\),
the term \(p=0\) gives a nonzero section and the same equality follows.
This is the equality \(r=r^-\) predicted in the remark ending
\cite[Section~3]{BrunebarbeCampana2015}.

For \textup{(ii)}, the forward vanishing is the rational-curve argument
recalled by Brunebarbe and Campana in the same section. Conversely, the
stated vanishing gives
\(H^{2,0}(X)=H^0(X,\Omega_X^2)=0\). Kodaira's projectivity criterion
then makes \(X\) projective: the real degree-two cohomology is of type
\((1,1)\), so a rational class sufficiently close to a K\"ahler class
is still K\"ahler, and clearing denominators gives an integral K\"ahler
class. This is the projectivity reduction in
\cite[Section~3]{BrunebarbeCampana2015}. Part \textup{(i)} now applies.
Its maximum is zero because every term with \(p>0\) vanishes, so the
MRC base is a point and \(X\) is rationally connected.
\end{proof}

We next use ordinary cotangent invariants, with no boundary. If
\(\mathcal L\) is a coherent rank-one subsheaf of a vector bundle on a
smooth projective variety, its double dual is a line bundle; put
\(\kappa(\mathcal L)=\kappa(\mathcal L^{**})\). Define
\[
 \omega(X)=
 \max_{\substack{m>0\\
 \mathcal L\subset(\Omega_X^1)^{\otimes m}\\
 \operatorname{rank}\mathcal L=1}}
 \kappa(\mathcal L),
\]
where \(m\) is an integer and the value is \(-\infty\) if no such sheaf
has nonnegative Iitaka dimension. The inclusion of \(\mathcal L\) extends
to its double dual: away from codimension two these sheaves agree, and a
map into a locally free sheaf extends across that subset. Taking tensor
powers then realizes the systems \(H^0(X,(\mathcal L^{**})^{\otimes a})\),
for integers \(a>0\), inside higher cotangent tensor powers. Thus allowing all \(m\) makes
this equivalent to Campana's definition using the dimensions of the
rational maps defined by \(H^0(X,\mathcal L)\). For positive-dimensional \(X\),
also put
\[
 \kappa^+(X)=
 \max_{\substack{p>0\\0\ne\mathcal F\subset\Omega_X^p}}
 \kappa(\det\mathcal F).
\]
Here \(p\) is an integer, \(\mathcal F\) is coherent, and
\(\det\mathcal F=(\bigwedge^{\rk\mathcal F}\mathcal F)^{**}\). These are the
invariants used in Campana's Appendix~A to Taji \cite{Taji2014}.

\begin{corollary}[Cotangent invariants and the MRC base]
\label{cor:cotangent-mrc}
Let \(X\) be a smooth connected projective complex variety. Then
\[
 \omega(X)=\omega(R(X))=
 \begin{cases}
 -\infty,&\text{if }R(X)\text{ is a point},\\
 \kappa(R(X)),&\text{if }\dim R(X)>0.
 \end{cases}
\]
In particular, if \(X\) is positive-dimensional and non-uniruled, then
\(\omega(X)=\kappa(X)\). If \(X\) is positive-dimensional and
\(\kappa(X)\geq0\), then
\[
 \kappa^+(X)=\kappa(X).
\]
\end{corollary}
\begin{proof}
Campana's birationally invariant fibration properties
\cite[Appendix~A, property~4]{Taji2014} give
\(\omega(X)=\omega(R(X))\): the general MRC fiber is rationally
connected and has \(\omega=-\infty\). Indeed any rank-one subsheaf with
nonnegative Iitaka dimension would give a section of a positive tensor
power, contrary to Corollary~\ref{cor:uniruled-mumford}; a point has
\(\omega=-\infty\) by the empty maximum. Birational invariance permits
resolving the MRC map. If \(R(X)\) is a point, the displayed formula
therefore follows without changing our convention \(\kappa(\mathrm{pt})=0\).

Suppose that \(R=R(X)\) has positive dimension. It is non-uniruled, so
\(\kappa(R)\geq0\) by Corollary~\ref{cor:uniruled-mumford}.
To apply Campana's abundance deduction
\cite[Conjecture~A.2 and Remark~A.3]{Taji2014}, take a resolved Iitaka
fibration \(R'\to Z\), with \(R'\) and \(Z\) smooth and projective,
\(R'\) birational to \(R\), and \(\dim Z=\kappa(R)\). Its smooth
general fiber \(F\) has \(\kappa(F)=0\). If \(F\) is
positive-dimensional, Corollary~\ref{cor:smooth-good-model} and
Lemma~\ref{lem:semiample-zero} give a normal birational good model
\(V_F\) with \(K_{V_F}\sim_{\Q}0\). In particular its canonical
divisor is numerically trivial on its smooth locus. Campana's criterion
\cite[Appendix~A, Examples~A.1(2)]{Taji2014} says that the existence of
such a normal birational model implies \(\omega(F)=0\); it does not
require \(K_F\) itself to be trivial. Property~4 of the same appendix
and birational invariance now give
\(\omega(R)=\omega(R')\leq\dim Z=\kappa(R)\).
If \(F\) is a point, this upper bound is simply
\(\omega(R)\leq\dim R=\kappa(R)\). The reverse inequality
\(\omega(R)\geq\kappa(R)\) is one of the defining cotangent inequalities
recorded in the same appendix. This proves the piecewise formula.
A non-uniruled \(X\) has \(R(X)\) birational to \(X\), giving the
first stated specialization.

Campana also records
\(\omega(X)\geq\kappa^+(X)\geq\kappa(X)\)
\cite[Appendix~A, property~2]{Taji2014}. When \(\kappa(X)\geq0\),
Corollary~\ref{cor:uniruled-mumford} makes \(X\) non-uniruled, and the
two outer terms are equal by the formula just proved. Thus
\(\kappa^+(X)=\kappa(X)\). This is the abundance consequence recalled
by Taji \cite[Theorem~1.3]{Taji2014}, who attributes it to Campana
\cite[Proposition~3.10]{Campana1995}.
\end{proof}

\begin{corollary}[Finite fundamental group in Kodaira dimension zero]
\label{cor:finite-pi-one}
Let \(X\) be a positive-dimensional smooth connected projective complex
variety. If \(\kappa(X)=0\) and \(\chi(X,\OO_X)\ne0\), then its ordinary
topological fundamental group \(\pi_1(X)\) is finite.
\end{corollary}
\begin{proof}
Corollary~\ref{cor:cotangent-mrc} gives \(\kappa^+(X)=0\).
Campana's criterion \cite[Corollary~5.3]{Campana1995}, in the form
recalled by Taji \cite[Theorem~1.6]{Taji2014}, gives the conclusion.
\end{proof}

\subsection{Quasi-projective covers}

Here a quasi-projective cover means a complex analytic covering space
biholomorphic to a quasi-projective variety; its deck transformations
need not be algebraic. The following application describes such covers
in terms of bundles over abelian varieties.

\begin{corollary}[Quasi-projective covering spaces]
\label{cor:quasi-projective-covers}
Let \(X\) be a connected normal projective complex variety.
\begin{enumerate}[label=\textup{(\roman*)}]
\item Its universal cover is biholomorphic to a quasi-projective variety
if and only if there is a finite \'etale Galois cover \(X'\to X\)
and a locally trivial holomorphic fiber bundle \(X'\to A\), where
\(A\) is an abelian variety and the fiber is simply connected and projective.
\item If \(\widetilde X\to X\) is an infinite \'etale Galois cover
and \(\widetilde X\) is biholomorphic to a quasi-projective variety,
there are a finite \'etale Galois cover \(X'\to X\), a locally trivial
holomorphic fiber bundle \(\alpha:X'\to A\) over an abelian variety,
and an \'etale cover \(\widetilde A\to A\) with no positive-dimensional
compact analytic subvarieties, such that
\[
 \widetilde X\simeq X'\times_A\widetilde A
\]
as covering spaces over \(X\).
\end{enumerate}
\end{corollary}
\begin{proof}
Theorem~\ref{thm:main} with zero boundary on a smooth projective complex
variety gives precisely Conjecture~1.2 of
Claudon--H\"oring--Koll\'ar \cite{ClaudonHoeringKollar2013}.
Their Theorem~1.1 and Corollary~1.5 therefore give \textup{(i)} and
\textup{(ii)}, respectively.
\end{proof}

The simply connected fiber condition belongs to \textup{(i)}, not to
the general-cover statement \textup{(ii)}. Local triviality is holomorphic;
no global product after a finite cover or compact K\"ahler extension is
asserted.
